\section{Computational certificates and reproducibility}
\label{cert:section}

The arguments use two different finite certificates.  The W12--W16 chain
certifies exact Fourier identities, analytic time-profile bounds and rational
inequalities for the whole-time Picard and endpoint arguments.  It does not
numerically integrate a Navier--Stokes trajectory.  The supplementary N3
certificate encloses a finite-dimensional Galerkin reference by outward
interval arithmetic and combines that enclosure with an analytic tail.  The
full-PDE conclusions in both cases also require the written analytic estimates
and the cutoff-uniform continuation argument.  A successful program exit alone
does not establish these analytic implications.

\subsection{Arithmetic and the mathematical objects certified}

The decisive W12--W16 inequalities are evaluated with integers and exact
rational numbers.  Square roots and exponential or logarithmic expressions
needed in the estimates are replaced by rational enclosures with explicit
remainder bounds.  Decimal conversions in the output are for orientation;
they are not used to decide the strict certificate inequalities.  In
particular, the smallest positive rectangle-invariance slack in the W15
data-only audit is approximately
\[
  1.1507455879810338\times 10^{-19}.
\]
This displayed decimal is not the proof of positivity.  The audit computes
each slack as an exact rational number and tests it strictly against zero.
Rounding the displayed slacks to a few decimal places would therefore destroy
essential information.  The rational coefficients and the complete inequality
rows are retained in the accompanying JSON files.

The W15 finite checks verify the full 17-block memory terms, invariance of the
rectangle, entry of the actual heat-seeded third iterate, and the estimate
relative to the fixed point using the sum of the third-iterate entry error and
the rectangle radius.  The geometric-rate and solution-stability checks use
positive rational majorants.  The W16 checks certify the separate endpoint
regions and their retained-history budgets.  The replay wrapper also binds
all nine field caps used in W16 to the freshly reconstructed W12 Fourier
bounds, and checks the full formal twelve-monomial residual expansion.  This
last check supplements, rather than replaces, the Fourier reconstruction.

For N3, an interval is stored as an integer pair representing
\([\ell/2^p,r/2^p]\).  The implementation uses outward integer rounding,
integer square-root bounds and rational series bounds for exponentials.
Binary floating-point arithmetic does not determine the passing gate.  The
reproduced reference uses \(p=224\), Taylor order \(16\), step \(1/128\), and
\(1024\) steps to time \(8\).  The analytic continuation of the bounds from
time \(8\) to infinity uses the finite ODE energy decay.  The reference
contains ten basis coordinates and \(108\) signed Fourier modes.  Its complete
exterior residual, rather than only the retained modes, enters the PDE error
estimate.  The accompanying rational margin check verifies, among other
inequalities, \(\kappa<49/100\).  The alternative
\(384\)-bit, order-\(32\) computation retained in the larger research archive
is not included in this selected bundle and was not rerun in the
reproduction reported here.

\subsection{Source-to-output map}

Paths in this appendix are relative to the extracted companion
\path{Code.zip}. The original research files are preserved under
\path{proof_sources/W_research/}.  The main analytic dependencies are
\path{07_nonlinear_heat/heat_space_fixed_point.md},
\path{09_geometric_growth/endpoint_critical_gate_audit.md}, and the W12--W16
proof notes identified in \path{SOURCE_MAP.json}.  For N3, the written bridge
from the finite reference to the PDE is in
\path{01_N3_W/n3_pde_scalar_audit.md}; the finite enclosures and tail are
described in \path{01_N3_W/n3_numerical_certificate.md}.

Table~\ref{tab:w-certificates} lists all fifteen distinct W programs rerun in
dependency order.  A directory label such as W15 means the corresponding
directory below \path{proof_sources/W_research/}: W12 is
\path{12_second_correction/}, W13 is \path{13_iteration_convergence/},
W14 is \path{14_unit_W_picard/}, W15 is \path{15_perturbed_picard/}, and
W16 is \path{16_enlarged_endpoint/}.  For every listed program, its result
filename is obtained by replacing \path{.py} with \path{_result.json}, with
one exception: \path{refined_W_block_gate_check.py} writes
\path{refined_W_block_gate_result.json}.  Thus the table, together with this
rule, specifies each source and output exactly. Archived comparison outputs
remain beside their original source programs. The replay wrapper places
copies of the programs and freshly generated dependencies together in
\path{working_copy/} under the chosen output directory. Individual programs
are not intended to be run directly from the separated original-stage
directories. Duplicate fresh outputs are omitted from this compact bundle;
the release replay receipt is retained. The source map identifies all 61
selected mathematical dependencies and their unchanged original bytes.
Historical headings in the original notes are retained as source evidence;
the consolidated manuscript states the present results.

{\small
\begin{longtable}{@{}p{0.07\textwidth}>{\raggedright\arraybackslash}p{0.53\textwidth}>{\raggedright\arraybackslash}p{0.33\textwidth}@{}}
\caption{Exact-arithmetic W certificate chain.}\label{tab:w-certificates}\\
\toprule
Stage & Source program & Certified component \\
\midrule
\endfirsthead
\toprule
Stage & Source program & Certified component \\
\midrule
\endhead
\bottomrule
\endfoot
W12 & \path{second_correction_fourier_check.py} & Exact second-correction Fourier fields and caps. \\
W12 & \path{second_corrected_gate_static.py} & Analytic second-reference gate. \\
W13 & \path{third_heat_defect_static.py} & Complete third heat-defect profiles. \\
W14 & \path{fourth_heat_defect_static.py} & Complete fourth heat-defect bounds. \\
W14 & \path{refined_W_block_gate_check.py} & Whole-history block rectangle. \\
W14 & \path{refined_iteration_adversarial_static.py} & Separate data-only rectangle audit. \\
W14 & \path{W_picard_rate_static.py} & Unit-W geometric-rate majorants. \\
W15 & \path{perturbation_block_gate.py} & Perturbed invariant rectangle. \\
W15 & \path{perturbed_picard_prefix_static.py} & Actual heat-seeded finite-prefix entry. \\
W15 & \path{perturbed_stability_rate_static.py} & Rates, solution budget and stability. \\
W15 & \path{perturbed_picard_adversarial_static.py} & Separate data-only audit of entry and rates. \\
W16 & \path{enlarged_endpoint_gate.py} & Five separate endpoint regularity regions. \\
W16 & \path{enlarged_budget_static.py} & Retained-history solution budgets. \\
W16 & \path{enlarged_endpoint_adversarial_static.py} & Separate data-only endpoint/budget audit. \\
W16 & \path{enlarged_class_inclusion_static.py} & Inclusion of the earlier analytic classes. \\
\end{longtable}
}

The N3 input is \path{n3_exact_algebra_intermediate.json}, which records the
exact basis, ODE coefficients and exterior residual. The original analytic
notes are in \path{proof_sources/W_research/01_N3_W/}. The executable
programs and this algebra input are in \path{checks/N3/}, with unchanged
file bytes; their original locations are recorded in \path{SOURCE_MAP.json}.
Archived comparison outputs are in \path{checks/N3/archived/}. The runner
generates fresh outputs in its own directory; duplicate fresh outputs are
omitted here while the release receipt is retained. Table~\ref{tab:n3-certificates}
gives the output map;
\path{n3_interval_core.py} is the arithmetic and integration module used by
the finite validator.

{\small
\begin{longtable}{@{}>{\raggedright\arraybackslash}p{0.34\textwidth}>{\raggedright\arraybackslash}p{0.39\textwidth}>{\raggedright\arraybackslash}p{0.20\textwidth}@{}}
\caption{N3 reproduction map.}\label{tab:n3-certificates}\\
\toprule
Source program & Output & Role \\
\midrule
\endfirsthead
\toprule
Source program & Output & Role \\
\midrule
\endhead
\bottomrule
\endfoot
\path{n3_algebra_audit.py} & \path{n3_algebra_audit_result.json} & Exact Fourier reconstruction. \\
\path{n3_validated_run.py} & \path{n3_validated_result.json} & Interval evolution, diagnostics and analytic tail. \\
\path{n3_check_margins.py} & \path{n3_rational_margin_check.json} & Separate rational gate and rounded bounds. \\
\path{run_replay.py} & \path{replay_receipt.json} & Commands, source hashes and archived comparisons. \\
\end{longtable}
}

\subsection{Reproduction record and commands}

The version 1.0 W replay ran on 16 October 2026, from 23:55:54 to 23:56:21 UTC
(17 October, 05:25:54--05:26:21 Indian Standard Time). All fifteen programs
returned success, and all fifteen result files were byte-for-byte identical
to their archived counterparts.  Inputs were copied without patches into a
new working directory; result files were generated afresh in dependency
order.  The original selected source files were unchanged.  The replayable
commands, result hashes, assertion counts where supplied, and dependency
checks are recorded in \path{checks/W/replay_receipt.json}.

The version 1.0 N3 replay ran on 16 October 2026, from 23:56:08 to 23:57:02 UTC
(17 October, 05:26:08--05:27:02 Indian Standard Time). The algebra audit,
finite validator and rational margin check all returned success.  The
mathematical fields of the finite result agreed exactly with the archived
JSON after excluding the runtime field \path{elapsed_seconds}; the algebra
and rational-margin JSON data also agreed exactly.  This is a statement of
parsed-data equality, not a claim that the entire N3 result file was
byte-identical.  Source hashes and comparisons appear in
\path{checks/N3/replay_receipt.json}.  Both reproductions used Python
3.12.14 and standard-library dependencies.

From the extracted bundle root, the W chain can be repeated into an absent or
empty output directory with:
\begin{verbatim}
python3 checks/replay_W_chain.py \
  --source proof_sources/W_research --out replay_new_W
\end{verbatim}
For N3, first make a separate copy of \path{checks/N3/}; from that copy run:
\begin{verbatim}
python3 run_replay.py
\end{verbatim}
The N3 runner writes outputs within its own directory. A separate copy
therefore preserves the packaged reproduction record. The W wrapper uses
portable defaults and records logical relative locations; the scientific
checker programs and archived certificates remain unchanged. Do not disable Python
assertions with optimization flags.  A later execution naturally changes
timestamps and wall-clock durations, so its entire receipt need not match
the package manifest. Receipts record log names for a repeated run; raw
execution logs and duplicate working trees are not distributed here.
The supplied manuscript sources, selected proof notes, programs, original
comparison certificates and retained version 1.0 receipts form the
reproducibility record. Version 1.1 changes exposition and references;
the scientific sources and receipts are unchanged, and no new scientific
replay is claimed for this revision.

\subsection{Evidence boundaries}

These reproductions establish that the supplied programs, on the recorded
runtime, reconstruct the archived passing certificates.  They do not provide
a separate implementation of the complete Fourier and interval algorithms.
The data-only audit programs check important identities and inequalities
separately, but do not constitute a formalization of every analytic lemma.
The proof still depends on the correctness of the mathematical reduction,
the implementations of the finite checks and enclosures, Python integer and
rational arithmetic, and the ordinary execution environment.  Standard PDE
local existence, compactness, uniqueness and continuation arguments also
remain part of the written proof.

The package manifest and SHA-256 digests serve file identification and
integrity checking.  They establish neither mathematical correctness nor an
independently authenticated historical date.  No formal proof-assistant
verification or external human validation is represented by a passing
receipt.  In particular, the W16 regions do not acquire the W15 ordinary
heat-seeded Picard-convergence conclusion merely because both certificate
chains pass, and the fixed-W N3 bound does not extend unchanged to all
perturbation regions.  Neither numerical reproduction nor source hashing
establishes novelty, priority, or regularity for arbitrary smooth
three-dimensional data.
